% Copyright 2006 by Till Tantau
%
% This file may be distributed and/or modified
%
% 1. under the LaTeX Project Public License and/or
% 2. under the GNU Free Documentation License.
%
% See the file doc/generic/pgf/licenses/LICENSE for more details.


\section{Tutorial: Diagrams as Simple Graphs}

In this tutorial we have a look at how graphs and matrices can be used to
typeset a diagram.

Ilka, who just got tenure for her professorship on Old and Lovable Programming
Languages, has recently dug up a technical report entitled \emph{The
Programming Language Pascal} in the dusty cellar of the library of her
university. Having been created in the good old times using pens and rules, it
looks like this%
\footnote{The shown diagram was not scanned, but rather typeset using
\tikzname. The jittering lines were created using the |random steps|
decoration.}:

{
  \tikzset{
    nonterminal/.style={
      % The shape:
      rectangle,
      % The size:
      minimum size=6mm,
      % The border:
      very thick,
      draw=red!50!black!50,         % 50% red and 50% black,
                                    % and that mixed with 50% white
      % The filling:
      top color=white,              % a shading that is white at the top...
      bottom color=red!50!black!20, % and something else at the bottom
      % Font
      font=\itshape
    },
    terminal/.style={
      % The shape:
      rounded rectangle,
      minimum size=6mm,
      % The rest
      very thick,draw=black!50,
      top color=white,bottom color=black!20,
      font=\ttfamily},
    skip loop/.style={to path={-- ++(0,#1) -| (\tikztotarget)}}
  }
  \tikzset{terminal/.append style={text height=1.5ex,text depth=.25ex}}
  \tikzset{nonterminal/.append style={text height=1.5ex,text depth=.25ex}}
  \pgfmathsetseed{1}
\medskip
\noindent\begin{tikzpicture}[
  >=latex,thick,
  /pgf/every decoration/.style={/tikz/sharp corners},
  fuzzy/.style={decorate,decoration={random steps,segment length=0.5mm,amplitude=0.15pt}},
  minimum size=6mm,line join=round,line cap=round,
  terminal/.style={rectangle,draw,fill=white,fuzzy,rounded corners=3mm},
  nonterminal/.style={rectangle,draw,fill=white,fuzzy},
  node distance=4mm]

  \ttfamily
  \begin{scope}[start chain,
                every node/.style={on chain},
                terminal/.append style={join=by {->,shorten >=-1pt,fuzzy,decoration={post length=4pt}}},
                nonterminal/.append style={join=by {->,shorten >=-1pt,fuzzy,decoration={post length=4pt}}},
                support/.style={coordinate,join=by fuzzy}]
    \node [support]             (start)        {};
    \node [nonterminal]                        {unsigned integer};
    \node [support]             (after ui)     {};
    \node [terminal]                           {.};
    \node [support]             (after dot)    {};
    \node [terminal]                           {digit};
    \node [support]             (after digit)  {};
    \node [support]             (skip)         {};
    \node [support]             (before E)     {};
    \node [terminal]                           {E};
    \node [support]             (after E)      {};
    \node [support,xshift=5mm]  (between)      {};
    \node [support,xshift=5mm]  (before last)  {};
    \node [nonterminal]                        {unsigned integer};
    \node [support]             (after last)   {};
    \node [coordinate,join=by ->] (end)          {};
  \end{scope}
  \node (plus)  [terminal,above=of between] {+};
  \node (minus) [terminal,below=of between] {-};

  \begin{scope}[->,decoration={post length=4pt},rounded corners=2mm,every path/.style=fuzzy]
    \draw (after ui)    -- +(0,.7)  -| (skip);
    \draw (after digit) -- +(0,-.7) -| (after dot);
    \draw (before E)    -- +(0,-1.2) -| (after last);
    \draw (after E)     |- (plus);
    \draw (plus)        -| (before last);
    \draw (after E)     |- (minus);
    \draw (minus)       -| (before last);
  \end{scope}
\end{tikzpicture}
\medskip

For her next lecture, Ilka decides to redo this diagram, but this time perhaps
a bit cleaner and perhaps also bit ``cooler''.

\medskip
\noindent\begin{tikzpicture}[point/.style={coordinate},>={Stealth[round]},thick,draw=black!50,
                    tip/.style={->,shorten >=1pt},every join/.style={rounded corners},
                    hv path/.style={to path={-| (\tikztotarget)}},
                    vh path/.style={to path={|- (\tikztotarget)}}]
  \matrix[column sep=4mm] {
    % First row:
    & & & & & & &  & & & & \node (plus) [terminal] {+};\\
    % Second row:
    \node (p1) [point]  {}; &    \node (ui1)   [nonterminal] {unsigned integer}; &
    \node (p2) [point]  {}; &    \node (dot)   [terminal]    {.};                &
    \node (p3) [point]  {}; &    \node (digit) [terminal]    {digit};            &
    \node (p4) [point]  {}; &    \node (p5)    [point]  {};                      &
    \node (p6) [point]  {}; &    \node (e)     [terminal]    {E};                &
    \node (p7) [point]  {}; &                                                    &
    \node (p8) [point]  {}; &    \node (ui2)   [nonterminal] {unsigned integer}; &
    \node (p9) [point]  {}; &    \node (p10)   [point]       {};\\
    % Third row:
    & & & & & & &  & & & & \node (minus)[terminal] {-};\\
  };

  { [start chain]
    \chainin (p1);
    \chainin (ui1)   [join=by tip];
    \chainin (p2)    [join];
    \chainin (dot)   [join=by tip];
    \chainin (p3)    [join];
    \chainin (digit) [join=by tip];
    \chainin (p4)    [join];
    { [start branch=digit loop]
      \chainin (p3) [join=by {skip loop=-6mm,tip}];
    }
    \chainin (p5)    [join,join=with p2 by {skip loop=6mm,tip}];
    \chainin (p6)    [join];
    \chainin (e)     [join=by tip];
    \chainin (p7)    [join];
    { [start branch=plus]
      \chainin (plus)  [join=by {vh path,tip}];
      \chainin (p8)    [join=by {hv path,tip}];
    }
    { [start branch=minus]
      \chainin (minus) [join=by {vh path,tip}];
      \chainin (p8)    [join=by {hv path,tip}];
    }
    \chainin (p8)    [join];
    \chainin (ui2)   [join=by tip];
    \chainin (p9)    [join,join=with p6 by {skip loop=-11mm,tip}];
    \chainin (p10)   [join=by tip];
  }
\end{tikzpicture}
}\medskip

Having read the previous tutorials, Ilka knows already how to set up the
environment for her diagram, namely using a |tikzpicture| environment. She
wonders which libraries she will need. She decides that she will postpone the
decision and add the necessary libraries as needed as she constructs the
picture.


\subsection{Styling the Nodes}

The bulk of this tutorial will be about arranging the nodes and connecting them
using chains, but let us start with setting up styles for the nodes.

There are two kinds of nodes in the diagram, namely what theoreticians like to
call \emph{terminals} and \emph{nonterminals}. For the terminals, Ilka decides
to use a black color, which visually shows that ``nothing needs to be done
about them''. The nonterminals, which still need to be ``processed'' further,
get a bit of red mixed in.

Ilka starts with the simpler nonterminals, as there are no rounded corners
involved. Naturally, she sets up a style:
%
\begin{codeexample}[preamble={\usetikzlibrary{positioning}}]
\begin{tikzpicture}[
    nonterminal/.style={
      % The shape:
      rectangle,
      % The size:
      minimum size=6mm,
      % The border:
      very thick,
      draw=red!50!black!50,         % 50% red and 50% black,
                                    % and that mixed with 50% white
      % The filling:
      top color=white,              % a shading that is white at the top...
      bottom color=red!50!black!20, % and something else at the bottom
      % Font
      font=\itshape
    }]
  \node [nonterminal] {unsigned integer};
\end{tikzpicture}
\end{codeexample}
%
Ilka is pretty proud of the use of the |minimum size| option: As the name
suggests, this option ensures that the node is at least 6mm by 6mm, but it will
expand in size as necessary to accommodate longer text. By giving this option
to all nodes, they will all have the same height of 6mm.

Styling the terminals is a bit more difficult because of the round corners.
Ilka has several options how she can achieve them. One way is to use the
|rounded corners| option. It gets a dimension as parameter and causes all
corners to be replaced by little arcs with the given dimension as radius. By
setting the radius to 3mm, she will get exactly what she needs: circles, when
the shapes are, indeed, exactly 6mm by 6mm and otherwise half circles on the
sides:
%
\begin{codeexample}[preamble={\usetikzlibrary{positioning}}]
\begin{tikzpicture}[node distance=5mm,
                    terminal/.style={
                      % The shape:
                      rectangle,minimum size=6mm,rounded corners=3mm,
                      % The rest
                      very thick,draw=black!50,
                      top color=white,bottom color=black!20,
                      font=\ttfamily}]
  \node (dot)   [terminal]                {.};
  \node (digit) [terminal,right=of dot]   {digit};
  \node (E)     [terminal,right=of digit] {E};
\end{tikzpicture}
\end{codeexample}

Another possibility is to use a shape that is specially made for typesetting
rectangles with arcs on the sides (she has to use the |shapes.misc| library to
use it). This shape gives Ilka much more control over the appearance. For
instance, she could have an arc only on the left side, but she will not need
this.
%
\begin{codeexample}[preamble={\usetikzlibrary{positioning,shapes.misc}}]
\begin{tikzpicture}[node distance=5mm,
                    terminal/.style={
                      % The shape:
                      rounded rectangle,
                      minimum size=6mm,
                      % The rest
                      very thick,draw=black!50,
                      top color=white,bottom color=black!20,
                      font=\ttfamily}]
  \node (dot)   [terminal]                {.};
  \node (digit) [terminal,right=of dot]   {digit};
  \node (E)     [terminal,right=of digit] {E};
\end{tikzpicture}
\end{codeexample}
%
At this point, she notices a problem. The baseline of the text in the nodes is
not aligned:
%
\begin{codeexample}[setup code,hidden]
\tikzset{
  terminal/.style={
    % The shape:
    rounded rectangle,
    minimum size=6mm,
    % The rest
    very thick,draw=black!50,
    top color=white,bottom color=black!20,
    font=\ttfamily},
}
\end{codeexample}
%
\begin{codeexample}[preamble={\usetikzlibrary{calc,positioning,shapes.misc}}]
\begin{tikzpicture}[node distance=5mm]
  \node (dot)   [terminal]                {.};
  \node (digit) [terminal,right=of dot]   {digit};
  \node (E)     [terminal,right=of digit] {E};

  \draw [help lines] let \p1 = (dot.base),
                         \p2 = (digit.base),
                         \p3 = (E.base)
                     in (-.5,\y1) -- (3.5,\y1)
                        (-.5,\y2) -- (3.5,\y2)
                        (-.5,\y3) -- (3.5,\y3);
\end{tikzpicture}
\end{codeexample}
%
\noindent (Ilka has moved the style definition to the preamble by saying
|\tikzset{terminal/.style=...}|, so that she can use it in all pictures.)

For the |digit| and the |E| the difference in the baselines is almost
imperceptible, but for the dot the problem is quite severe: It looks more like
a multiplication dot than a period.

Ilka toys with the idea of using the |base right=of...| option rather than
|right=of...| to align the nodes in such a way that the baselines are all on
the same line (the |base right| option places a node right of something so that
the baseline is right of the baseline of the other object). However, this does
not have the desired effect:
%
\begin{codeexample}[preamble={\usetikzlibrary{positioning,shapes.misc}}]
\begin{tikzpicture}[node distance=5mm]
  \node (dot)   [terminal]                     {.};
  \node (digit) [terminal,base right=of dot]   {digit};
  \node (E)     [terminal,base right=of digit] {E};
\end{tikzpicture}
\end{codeexample}
%
The nodes suddenly ``dance around''! There is no hope of changing the position
of text inside a node using anchors. Instead, Ilka must use a trick: The
problem of mismatching baselines is caused by the fact that |.| and |digit| and
|E| all have different heights and depth. If they all had the same, they would
all be positioned vertically in the same manner. So, all Ilka needs to do is to
use the |text height| and |text depth| options to explicitly specify a height
and depth for the nodes.
%
\begin{codeexample}[preamble={\usetikzlibrary{positioning,shapes.misc}}]
\begin{tikzpicture}[node distance=5mm,
                    text height=1.5ex,text depth=.25ex]
  \node (dot)   [terminal]                {.};
  \node (digit) [terminal,right=of dot]   {digit};
  \node (E)     [terminal,right=of digit] {E};
\end{tikzpicture}
\end{codeexample}


\subsection{Aligning the Nodes Using Positioning Options}

Ilka now has the ``styling'' of the nodes ready. The next problem is to place
them in the right places. There are several ways to do this. The most
straightforward is to simply explicitly place the nodes at certain coordinates
``calculated by hand''. For very simple graphics this is perfectly alright, but
it has several disadvantages:
%
\begin{enumerate}
    \item For more difficult graphics, the calculation may become
        complicated.
    \item Changing the text of the nodes may make it necessary to recalculate
        the coordinates.
    \item The source code of the graphic is not very clear since the
        relationships between the positions of the nodes are not made
        explicit.
\end{enumerate}

For these reasons, Ilka decides to try out different ways of arranging the
nodes on the page.

The first method is the use of \emph{positioning options}. To use them, you
need to load the |positioning| library. This gives you access to advanced
implementations of options like |above| or |left|, since you can now say
|above=of some node| in order to place a node above of |some node|, with the
borders separated by |node distance|.

Ilka can use this to draw the place the nodes in a long row:
%
\begin{codeexample}[setup code,hidden]
\tikzset{
  nonterminal/.style={
    % The shape:
    rectangle,
    % The size:
    minimum size=6mm,
    % The border:
    very thick,
    draw=red!50!black!50,         % 50% red and 50% black,
                                  % and that mixed with 50% white
    % The filling:
    top color=white,              % a shading that is white at the top...
    bottom color=red!50!black!20, % and something else at the bottom
    % Font
    font=\itshape,
  },
}
\tikzset{
  terminal/.append style={text height=1.5ex,text depth=.25ex},
  nonterminal/.append style={text height=1.5ex,text depth=.25ex},
}
\end{codeexample}
%
\begin{codeexample}[preamble={\usetikzlibrary{positioning,shapes.misc}}]
\begin{tikzpicture}[node distance=5mm and 5mm]
  \node (ui1)   [nonterminal]                     {unsigned integer};
  \node (dot)   [terminal,right=of ui1]           {.};
  \node (digit) [terminal,right=of dot]           {digit};
  \node (E)     [terminal,right=of digit]         {E};
  \node (plus)  [terminal,above right=of E]       {+};
  \node (minus) [terminal,below right=of E]       {-};
  \node (ui2)   [nonterminal,below right=of plus] {unsigned integer};
\end{tikzpicture}
\end{codeexample}

For the plus and minus nodes, Ilka is a bit startled by their placements.
Shouldn't they be more to the right? The reason they are placed in that manner
is the following: The |north east| anchor of the |E| node lies at the ``upper
start of the right arc'', which, a bit unfortunately in this case, happens to
be the top of the node. Likewise, the |south west| anchor of the |+| node is
actually at its bottom and, indeed, the horizontal and vertical distances
between the top of the |E| node and the bottom of the |+| node are both 5mm.

There are several ways of fixing this problem. The easiest way is to simply add
a little bit of horizontal shift by hand:
%
\begin{codeexample}[preamble={\usetikzlibrary{positioning,shapes.misc}}]
\begin{tikzpicture}[node distance=5mm and 5mm]
  \node (E)     [terminal]                                   {E};
  \node (plus)  [terminal,above right=of E,xshift=5mm]       {+};
  \node (minus) [terminal,below right=of E,xshift=5mm]       {-};
  \node (ui2)   [nonterminal,below right=of plus,xshift=5mm] {unsigned integer};
\end{tikzpicture}
\end{codeexample}

A second way is to revert back to the idea of using a normal rectangle for the
terminals, but with rounded corners. Since corner rounding does not affect
anchors, she gets the following result:
%
\begin{codeexample}[preamble={\usetikzlibrary{positioning,shapes.misc}}]
\begin{tikzpicture}[node distance=5mm and 5mm,terminal/.append style={rectangle,rounded corners=3mm}]
  \node (E)     [terminal]                        {E};
  \node (plus)  [terminal,above right=of E]       {+};
  \node (minus) [terminal,below right=of E]       {-};
  \node (ui2)   [nonterminal,below right=of plus] {unsigned integer};
\end{tikzpicture}
\end{codeexample}
%
A third way is to use matrices, which we will do later.

Now that the nodes have been placed, Ilka needs to add connections. Here, some
connections are more difficult than others. Consider for instance the
``repeat'' line around the |digit|. One way of describing this line is to say
``it starts a little to the right of |digit| than goes down and then goes to
the left and finally ends at a point a little to the left of |digit|''. Ilka
can put this into code as follows:
%
\begin{codeexample}[preamble={\usetikzlibrary{calc,positioning,shapes.misc}}]
\begin{tikzpicture}[node distance=5mm and 5mm]
  \node (dot)   [terminal]                        {.};
  \node (digit) [terminal,right=of dot]           {digit};
  \node (E)     [terminal,right=of digit]         {E};

  \path (dot)   edge[->] (digit)  % simple edges
        (digit) edge[->] (E);

  \draw [->]
     % start right of digit.east, that is, at the point that is the
     % linear combination of digit.east and the vector (2mm,0pt). We
     % use the ($ ... $) notation for computing linear combinations
     ($ (digit.east) + (2mm,0) $)
     % Now go down
     -- ++(0,-.5)
     % And back to the left of digit.west
     -| ($ (digit.west) - (2mm,0) $);
\end{tikzpicture}
\end{codeexample}

Since Ilka needs this ``go up/down then horizontally and then up/down to a
target'' several times, it seems sensible to define a special \emph{to-path}
for this. Whenever the |edge| command is used, it simply adds the current value
of |to path| to the path. So, Ilka can set up a style that contains the correct
path:
%
\begin{codeexample}[preamble={\usetikzlibrary{calc,positioning,shapes.misc}}]
\begin{tikzpicture}[node distance=5mm and 5mm,
    skip loop/.style={to path={-- ++(0,-.5) -| (\tikztotarget)}}]
  \node (dot)   [terminal]                        {.};
  \node (digit) [terminal,right=of dot]           {digit};
  \node (E)     [terminal,right=of digit]         {E};

  \path (dot)   edge[->]           (digit)  % simple edges
        (digit) edge[->]           (E)
        ($ (digit.east) + (2mm,0) $)
                edge[->,skip loop] ($ (digit.west) - (2mm,0) $);
\end{tikzpicture}
\end{codeexample}

Ilka can even go a step further and make her |skip loop| style parameterized.
For this, the skip loop's vertical offset is passed as parameter |#1|. Also, in
the following code Ilka specifies the start and targets differently, namely as
the positions that are ``in the middle between the nodes''.
%
\begin{codeexample}[preamble={\usetikzlibrary{calc,positioning,shapes.misc}}]
\begin{tikzpicture}[node distance=5mm and 5mm,
    skip loop/.style={to path={-- ++(0,#1) -| (\tikztotarget)}}]
  \node (dot)   [terminal]                        {.};
  \node (digit) [terminal,right=of dot]           {digit};
  \node (E)     [terminal,right=of digit]         {E};

  \path (dot)   edge[->]                (digit)  % simple edges
        (digit) edge[->]                (E)
        ($ (digit.east)!.5!(E.west) $)
                edge[->,skip loop=-5mm] ($ (digit.west)!.5!(dot.east) $);
\end{tikzpicture}
\end{codeexample}


\subsection{Aligning the Nodes Using Matrices}

Ilka is still bothered a bit by the placement of the plus and minus nodes.
Somehow, having to add an explicit |xshift| seems too much like cheating.

A perhaps better way of positioning the nodes is to use a \emph{matrix}. In
\tikzname\ matrices can be used to align quite arbitrary graphical objects in
rows and columns. The syntax is very similar to the use of arrays and tables in
\TeX\ (indeed, internally \TeX\ tables are used, but a lot of stuff is going on
additionally).

In Ilka's graphic, there will be three rows: One row containing only the plus
node, one row containing the main nodes and one row containing only the minus
node.
%
\begin{codeexample}[preamble={\usetikzlibrary{shapes.misc}}]
\begin{tikzpicture}
  \matrix[row sep=1mm,column sep=5mm] {
    % First row:
      & & & & \node [terminal] {+}; & \\
    % Second row:
    \node [nonterminal] {unsigned integer}; &
    \node [terminal]    {.};                &
    \node [terminal]    {digit};            &
    \node [terminal]    {E};                &
                                            &
    \node [nonterminal] {unsigned integer}; \\
    % Third row:
      & & & & \node [terminal] {-}; & \\
  };
\end{tikzpicture}
\end{codeexample}
%
That was easy! By toying around with the row and columns separations, Ilka can
achieve all sorts of pleasing arrangements of the nodes.

Ilka now faces the same connecting problem as before. This time, she has an
idea: She adds small nodes (they will be turned into coordinates later on and
be invisible) at all the places where she would like connections to start and
end.
%
\begin{codeexample}[preamble={\usetikzlibrary{shapes.misc}}]
\begin{tikzpicture}[point/.style={circle,inner sep=0pt,minimum size=2pt,fill=red},
                   skip loop/.style={to path={-- ++(0,#1) -| (\tikztotarget)}}]
  \matrix[row sep=1mm,column sep=2mm] {
    % First row:
    & & & & & & &  & & & & \node (plus) [terminal] {+};\\
    % Second row:
    \node (p1) [point]  {}; &    \node (ui1)   [nonterminal] {unsigned integer}; &
    \node (p2) [point]  {}; &    \node (dot)   [terminal]    {.};                &
    \node (p3) [point]  {}; &    \node (digit) [terminal]    {digit};            &
    \node (p4) [point]  {}; &    \node (p5)    [point]  {};                      &
    \node (p6) [point]  {}; &    \node (e)     [terminal]    {E};                &
    \node (p7) [point]  {}; &                                                    &
    \node (p8) [point]  {}; &    \node (ui2)   [nonterminal] {unsigned integer}; &
    \node (p9) [point]  {}; &    \node (p10)   [point]       {};\\
    % Third row:
    & & & & & & &  & & & & \node (minus)[terminal] {-};\\
  };

  \path (p4) edge [->,skip loop=-5mm] (p3)
        (p2) edge [->,skip loop=5mm]  (p6);
\end{tikzpicture}
\end{codeexample}
%
Now, it's only a small step to add all the missing edges.


\subsection{The Diagram as a Graph}

Matrices allow Ilka to align the nodes nicely, but the connections are not
quite perfect. The problem is that the code does not really reflect the paths
that underlie the diagram. For this, it seems natural enough to Ilka to use the
|graphs| library since, after all, connecting nodes by edges is exactly what
happens in a graph. The |graphs| library can both be used to connect nodes that
have already been created, but it can also be used to create nodes ``on the
fly'' and these processes can also be mixed.


\subsubsection{Connecting Already Positioned Nodes}

Ilka has already a fine method for positioning her nodes (using a |matrix|), so
all that she needs is an easy way of specifying the edges. For this, she uses
the |\graph| command (which is actually just a shorthand for |\path graph|). It
allows her to write down edges between them in a simple way (the macro
|\matrixcontent| contains exactly the matrix content from the previous example;
no need to repeat it here):
%
\begin{codeexample}[setup code,hidden]
\def\matrixcontent{
  % First row:
  \& \& \& \& \& \& \&  \& \& \& \& \node (plus) [terminal] {+};\\
  % Second row:
  \node (p1) [point]  {}; \&    \node (ui1)   [nonterminal] {unsigned integer}; \&
  \node (p2) [point]  {}; \&    \node (dot)   [terminal]    {.};                \&
  \node (p3) [point]  {}; \&    \node (digit) [terminal]    {digit};            \&
  \node (p4) [point]  {}; \&    \node (p5)    [point]  {};                      \&
  \node (p6) [point]  {}; \&    \node (e)     [terminal]    {E};                \&
  \node (p7) [point]  {}; \&                                                    \&
  \node (p8) [point]  {}; \&    \node (ui2)   [nonterminal] {unsigned integer}; \&
  \node (p9) [point]  {}; \&    \node (p10)   [point]       {};\\
  % Third row:
  \& \& \& \& \& \& \&  \& \& \& \& \node (minus)[terminal] {-};\\
}
\end{codeexample}
%
\begin{codeexample}[
    preamble={\usetikzlibrary{graphs,shapes.misc}},
    pre={\tikzset{ampersand replacement=\&,point/.style={coordinate}}},
]
\begin{tikzpicture}[skip loop/.style={to path={-- ++(0,#1) -| (\tikztotarget)}},
                    hv path/.style={to path={-| (\tikztotarget)}},
                    vh path/.style={to path={|- (\tikztotarget)}}]
  \matrix[row sep=1mm,column sep=2mm] { \matrixcontent };

  \graph {
    (p1) -> (ui1) -- (p2) -> (dot) -- (p3) -> (digit) -- (p4)
         -- (p5)  -- (p6) -> (e) -- (p7) -- (p8) -> (ui2) -- (p9) -> (p10);
    (p4) ->[skip loop=-5mm]  (p3);
    (p2) ->[skip loop=5mm]   (p5);
    (p6) ->[skip loop=-11mm] (p9);
    (p7) ->[vh path]         (plus)  -> [hv path] (p8);
    (p7) ->[vh path]         (minus) -> [hv path] (p8);
  };
\end{tikzpicture}
\end{codeexample}

This is already pretty near to the desired result, just a few ``finishing
touches'' are needed to style the edges more nicely.

However, Ilka does not have the feeling that the |graph| command is all that
hot in the example. It certainly does cut down on the number of characters she
has to write, but the overall graph structure is not that much clear -- it is
still mainly a list of paths through the graph. It would be nice to specify
that, say, there the path from |(p7)| sort of splits to |(plus)| and |(minus)|
and then merges once more at |(p8)|. Also, all these parentheses are bit hard
to type.

It turns out that edges from a node to a whole group of nodes are quite easy to
specify, as shown in the next example. Additionally, by using the
|use existing nodes| option, Ilka can also leave out all the parentheses
(again, some options have been moved outside to keep the examples shorter):
%
\begin{codeexample}[
    preamble={\usetikzlibrary{arrows.meta,graphs,shapes.misc}},
    pre={\tikzset{
    ampersand replacement=\&,
    point/.style={coordinate},
    skip loop/.style={to path={-- ++(0,##1) -| (\tikztotarget)}},
    hv path/.style={to path={-| (\tikztotarget)}},
    vh path/.style={to path={|- (\tikztotarget)}},
}},
]
\begin{tikzpicture}[>={Stealth[round]},thick,black!50,text=black,
                    every new ->/.style={shorten >=1pt},
                    graphs/every graph/.style={edges=rounded corners}]
  \matrix[column sep=4mm] { \matrixcontent };

  \graph [use existing nodes] {
    p1 -> ui1 -- p2 -> dot -- p3 -> digit -- p4 -- p5  -- p6 -> e -- p7 -- p8 -> ui2 -- p9 -> p10;
    p4 ->[skip loop=-5mm]  p3;
    p2 ->[skip loop=5mm]   p5;
    p6 ->[skip loop=-11mm] p9;
    p7 ->[vh path] { plus, minus } -> [hv path] p8;
  };
\end{tikzpicture}
\end{codeexample}


\subsubsection{Creating Nodes Using the Graph Command}

Ilka has heard that the |graph| command is also supposed to make it easy to
create nodes, not only to connect them. This is, indeed, correct: When the
|use existing nodes| option is not used and when a node name is not surrounded
by parentheses, then \tikzname\ will actually create a node whose name and text
is the node name:
%
\begin{codeexample}[preamble={\usetikzlibrary{graphs}}]
\tikz \graph [grow right=2cm] { unsigned integer -> d -> digit -> E };
\end{codeexample}
%
Not quite perfect, but we are getting somewhere. First, let us change the
positioning algorithm by saying |grow right sep|, which causes new nodes to be
placed to the right of the previous nodes with a certain fixed separation
(|1em| by default). Second, we add some options to make the node ``look nice''.
Third, note the funny |d| node above: Ilka tried writing just |.| there first,
but got some error messages. The reason is that a node cannot be called |.| in
\tikzname, so she had to choose a different name -- which is not good, since
she wants a dot to be shown! The trick is to put the dot in quotation marks,
this allows you to use ``quite arbitrary text'' as a node name:
%
\begin{codeexample}[preamble={\usetikzlibrary{graphs,shapes.misc}}]
\tikz \graph [grow right sep] {
  unsigned integer[nonterminal] -> "."[terminal] -> digit[terminal] -> E[terminal]
};
\end{codeexample}
%
Now comes the fork to the plus and minus signs. Here, Ilka can use the grouping
mechanism of the |graph| command to create a split:
%
\begin{codeexample}[preamble={\usetikzlibrary{graphs,shapes.misc}}]
\tikz \graph [grow right sep] {
  unsigned integer  [nonterminal] ->
  "."               [terminal] ->
  digit             [terminal] ->
  E                 [terminal] ->
  {
    "+"             [terminal],
    ""              [coordinate],
    "-"             [terminal]
  } ->
  ui2/unsigned integer [nonterminal]
};
\end{codeexample}
%
Let us see, what is happening here. We want two |unsigned integer| nodes, but
if we just were to use this text twice, then \tikzname\ would have noticed that
the same name was used already in the current graph and, being smart (actually
too smart in this case), would have created an edge back to the already-created
node. Thus, a fresh name is needed here. However, Ilka also cannot just write
|unsigned integer2|, because she wants the original text to be shown, after
all! The trick is to use a slash inside the node name: In order to ``render''
the node, the text following the slash is used instead of the node name, which
is the text before the slash. Alternatively, the |as| option can be used, which
also allows you to specify how a node should be rendered.

It turns out that Ilka does not need to invent a name like |ui2| for a node
that she will not reference again anyway. In this case, she can just leave out
the name (write nothing before |/|), which always stands for a ``fresh,
anonymous'' node name.

Next, Ilka needs to add some coordinates in between of some nodes where the
back-loops should got and she needs to shift the nodes a bit:
%
\begin{codeexample}[
    preamble={\usetikzlibrary{arrows.meta,graphs,shapes.misc}},
    pre={\tikzset{
    skip loop/.style={to path={-- ++(0,##1) -| (\tikztotarget)}},
    hv path/.style={to path={-| (\tikztotarget)}},
    vh path/.style={to path={|- (\tikztotarget)}},
}},
]
\begin{tikzpicture}[>={Stealth[round]}, thick, black!50, text=black,
                    every new ->/.style={shorten >=1pt},
                    graphs/every graph/.style={edges=rounded corners}]
  \graph [grow right sep, branch down=7mm] {
    /                  [coordinate] ->
    unsigned integer   [nonterminal] --
    p1                 [coordinate] ->
    "."                [terminal] --
    p2                 [coordinate] ->
    digit              [terminal] --
    p3                 [coordinate] --
    p4                 [coordinate] --
    p5                 [coordinate] ->
    E                  [terminal] --
    q1                 [coordinate] ->[vh path]
    { [nodes={yshift=7mm}]
      "+"                [terminal],
      q2/                [coordinate],
      "-"                [terminal]
    } -> [hv path]
    q3                 [coordinate] --
    /unsigned integer  [nonterminal] --
    p6                 [coordinate] ->
    /                  [coordinate];

    p1 ->[skip loop=5mm]   p4;
    p3 ->[skip loop=-5mm]  p2;
    p5 ->[skip loop=-11mm] p6;
  };
\end{tikzpicture}
\end{codeexample}

All that remains to be done is to somehow get rid of the strange curves between
the |E| and the unsigned integer. They are caused by \tikzname's attempt at
creating an edge that first goes vertical and then horizontal but is actually
just horizontal. Additionally, the edge should not really be pointed; but it
seems difficult to get rid of this since the \emph{other} edges from |q1|,
namely to |plus| and |minus| should be pointed.

It turns out that there is a nice way of solving this problem: You can specify
that a graph is |simple|. This means that there can be at most one edge between
any two nodes. Now, if you specify an edge twice, the options of the second
specification ``win''. Thus, by adding two more lines that ``correct'' these
edges, we get the final diagram with its complete code:
%
\begin{codeexample}[preamble={\usetikzlibrary{arrows.meta,graphs,shapes.misc}}]
\tikz [>={Stealth[round]}, black!50, text=black, thick,
       every new ->/.style          = {shorten >=1pt},
       graphs/every graph/.style    = {edges=rounded corners},
       skip loop/.style             = {to path={-- ++(0,#1) -| (\tikztotarget)}},
       hv path/.style               = {to path={-| (\tikztotarget)}},
       vh path/.style               = {to path={|- (\tikztotarget)}},
       nonterminal/.style           = {
         rectangle, minimum size=6mm, very thick, draw=red!50!black!50, top color=white,
         bottom color=red!50!black!20, font=\itshape, text height=1.5ex,text depth=.25ex},
       terminal/.style              = {
         rounded rectangle,  minimum size=6mm, very thick, draw=black!50, top color=white,
         bottom color=black!20, font=\ttfamily, text height=1.5ex, text depth=.25ex},
       shape                        = coordinate
       ]
  \graph [grow right sep, branch down=7mm, simple] {
    / -> unsigned integer[nonterminal] -- p1 -> "." [terminal] -- p2 -> digit[terminal] --
    p3 -- p4 -- p5 -> E[terminal] -- q1 ->[vh path]
    {[nodes={yshift=7mm}]
      "+"[terminal], q2, "-"[terminal]
    } -> [hv path]
    q3 -- /unsigned integer [nonterminal] -- p6 -> /;

    p1 ->[skip loop=5mm]   p4;
    p3 ->[skip loop=-5mm]  p2;
    p5 ->[skip loop=-11mm] p6;

    q1 -- q2 -- q3;  % make these edges plain
  };
\end{codeexample}


%% TODOsp: a commented subsection
% \subsection{Using Chains}
%
% Matrices allow Ilka to align the nodes nicely, but the connections are
% not quite perfect. The problem is that the code does not really
% reflect the paths that underlie the diagram.
%
%
% For this reason, Ilka decides to try out \emph{chains} by including
% the |chain| library. Basically, a chain is just a sequence of
% (usually) connected nodes. The nodes can already have been constructed
% or they can be constructed as the chain is constructed (or these
% processes can be mixed).
%
% \subsubsection{Creating a Simple Chain}
%
%
% Ilka starts with creating a chain from scratch. For this, she starts a
% chain using the |start chain| option in a scope. Then, inside the
% scope, she uses the |on chain| option on nodes to add them to the
% chain.
% \begin{codeexample}[]
% \begin{tikzpicture}[start chain,node distance=5mm]
%   \node [on chain,nonterminal]  {unsigned integer};
%   \node [on chain,terminal]     {.};
%   \node [on chain,terminal]     {digit};
%   \node [on chain,terminal]     {E};
%   \node [on chain,nonterminal]  {unsigned integer};
% \end{tikzpicture}
% \end{codeexample}
% (Ilka will add the plus and minus nodes later.)
%
% As can be seen, the nodes of a chain are placed in a row. This can be
% changed, for instance by saying |start chain=going below| we get a
% chain where each node is below the previous one.
%
% The next step is to \emph{join} the nodes of the chain. For this, we
% add the |join| option to each node. This joins the node with the
% previous node (for the first node nothing happens).
% \begin{codeexample}[]
% \begin{tikzpicture}[start chain,node distance=5mm]
%   \node [on chain,join,nonterminal]  {unsigned integer};
%   \node [on chain,join,terminal]     {.};
%   \node [on chain,join,terminal]     {digit};
%   \node [on chain,join,terminal]     {E};
%   \node [on chain,join,nonterminal]  {unsigned integer};
% \end{tikzpicture}
% \end{codeexample}
% In order to get a arrow tip, we redefine the |every join| style. Also,
% we move the |join| and |on chain| options to the |every node|
% style so that we do not have to repeat them so often.
% \begin{codeexample}[]
% \begin{tikzpicture}[start chain,node distance=5mm, every node/.style={on chain,join}, every join/.style={->}]
%   \node [nonterminal]  {unsigned integer};
%   \node [terminal]     {.};
%   \node [terminal]     {digit};
%   \node [terminal]     {E};
%   \node [nonterminal]  {unsigned integer};
% \end{tikzpicture}
% \end{codeexample}
%
%
% \subsubsection{Branching and Joining a Chain}
%
% It is now time to add the plus and minus signs. They obviously
% \emph{branch off} the main chain. For this reason, we start a branch
% for them using the |start branch| option.
% \begin{codeexample}[]
% \begin{tikzpicture}[start chain,node distance=5mm, every node/.style={on chain,join}, every join/.style={->}]
%   \node [nonterminal]  {unsigned integer};
%   \node [terminal]     {.};
%   \node [terminal]     {digit};
%   \node [terminal]     {E};
%   \begin{scope}[start branch=plus]
%     \node (plus)  [terminal,on chain=going above right] {+};
%   \end{scope}
%   \begin{scope}[start branch=minus]
%     \node (minus) [terminal,on chain=going below right] {-};
%   \end{scope}
%   \node [nonterminal,join=with plus,join=with minus]  {unsigned integer};
% \end{tikzpicture}
% \end{codeexample}
%
% Let us see, what is going on here. First, the |start branch| begins a
% branch, starting with the node last created on the current chain,
% which is the |E| node in our case. This is implicitly also the first
% node on this branch. A branch is nothing different from a chain, which
% is why the plus node is put on this branch using the |on chain|
% option. However, this time we specify the placement of the node
% explicitly using |going |\meta{direction}. This causes the plus sign
% to be placed above and right of the |E| node. It is automatically
% joined to its predecessor on the branch by the implicit |join|
% option.
%
% When the first branch ends, only the plus node has been added and the
% current chain is the original chain once more and we are back to the
% |E| node. Now we start a new branch for the minus node. After this
% branch, the current chain ends at |E| node once more.
%
% Finally, the rightmost unsigned integer is added to the (main) chain,
% which is why it is joined correctly with the |E| node. The two
% additional |join| options get a special |with| parameter. This allows
% you to join a node with a node other than the predecessor on the
% chain. The  |with| should be followed by the name of a node.
%
% Since Ilka will need scopes more often in the following, she includes
% the |scopes| library. This allows her to replace |\begin{scope}|
%   simply by an opening brace and  |\end{scope}| by the corresponding
% closing brace. Also, in the following example we reference
% the nodes |plus| and |minus| using
% their automatic name: The $i$th node on a chain is called
% |chain-|\meta{i}. For a branch \meta{branch}, the $i$th node is called
% |chain/|\meta{branch}|-|\meta{i}. The \meta{i} can be replaced by
% |begin| and |end| to reference the first and (currently) last node on
% the chain.
%
% \begin{codeexample}[]
% \begin{tikzpicture}[start chain,node distance=5mm, every on chain/.style={join}, every join/.style={->}]
%   \node [on chain,nonterminal]  {unsigned integer};
%   \node [on chain,terminal]     {.};
%   \node [on chain,terminal]     {digit};
%   \node [on chain,terminal]     {E};
%   { [start branch=plus]
%     \node (plus)  [terminal,on chain=going above right] {+};
%   }
%   { [start branch=minus]
%     \node (minus) [terminal,on chain=going below right] {-};
%   }
%   \node [nonterminal,on chain,join=with chain/plus-end,join=with chain/minus-end]  {unsigned integer};
% \end{tikzpicture}
% \end{codeexample}
%
%
% The next step is to add intermediate coordinate nodes in the same
% manner as Ilka did for the matrix. For them, we change the |join|
% style slightly, namely for these nodes we do not want an arrow
% tip. This can be achieved either by (locally) changing the
% |every join| style or, which is what is done in the below example, by
% giving the desired style using |join=by ...|, where |...| is the style
% to be used for the join.
%
% \begin{codeexample}[]
% \begin{tikzpicture}[start chain,node distance=5mm and 2mm,
%                     every node/.style={on chain},
%                     nonterminal/.append style={join=by ->},
%                     terminal/.append style={join=by ->},
%                     point/.style={join=by -,circle,fill=red,minimum size=2pt,inner sep=0pt}]
%   \node [point] {};  \node [nonterminal] {unsigned integer};
%   \node [point] {};  \node [terminal]    {.};
%   \node [point] {};  \node [terminal]    {digit};
%   \node [point] {};  \node [point]       {};
%   \node [point] {};  \node [terminal]    {E};
%   \node [point] {};
%   { [node distance=5mm and 1cm] % local change in horizontal distance
%     { [start branch=plus]
%       \node (plus)  [terminal,on chain=going above right] {+};
%     }
%     { [start branch=minus]
%       \node (minus) [terminal,on chain=going below right] {-};
%     }
%     \node [point,below right=of plus,join=with chain/plus-end by ->,join=with chain/minus-end by ->] {};
%   }
%   \node [nonterminal] {unsigned integer};
%   \node [point]       {};
% \end{tikzpicture}
% \end{codeexample}
%
%
% \subsubsection{Chaining Together Already Positioned Nodes}
%
% The final step is to add the missing arrows. We can also use branches
% for them (even though we do not have to, but it is good practice and
% they exhibit the structure of the diagram in the code).
%
% Let us start with the repeat loop around the |digit|. This can be
% thought of as a branch that starts at the point after the digit and
% that ends at the point before the digit. However, we have already
% constructed the point before the digit! In such cases, it is possible
% to ``chain in'' an already positioned node, using the |\chainin|
% command. This command must be followed by a coordinate that contains a
% node name and optionally some options. The effect is that the named
% node is made part of the current chain.
%
% \begin{codeexample}[pre={\tikzset{node distance=5mm and 2mm,
%                     every node/.style={on chain},
%                     terminal/.append style={join=by ->},
%                     point/.style={join=by -,circle,fill=red,minimum size=2pt,inner sep=0pt}}}]
% \begin{tikzpicture}[start chain] % plus some styles that are not shown
%   \node                [point] {};
%   \node (before digit) [point] {};
%   \node                [terminal]    {digit};
%   \node                [point] {};
%   { [start branch=digit loop]
%     \chainin (before digit) [join=by {->,skip loop=-5mm}];
%   }
%   \node                [point] {};
% \end{tikzpicture}
% \end{codeexample}
%
%
% \subsubsection{Combined Use of Matrices and Chains}
%
% Ilka's final idea is to combine matrices and chains in the following
% manner: She will use a matrix to position the nodes. However, to show
% the logical ``flow structure'' inside the diagram, she will create
% chains and branches that show what is going on.
%
% Ilka starts with the matrix we had earlier, only with slightly adapted
% styles. Then she writes down the main chain and its branches:
%
% \begin{codeexample}[preamble={\usetikzlibrary{arrows.meta}}]
% \begin{tikzpicture}[point/.style={coordinate},>={Stealth[round]},thick,draw=black!50,
%                     tip/.style={->,shorten >=1pt},every join/.style={rounded corners},
%                     hv path/.style={to path={-| (\tikztotarget)}},
%                     vh path/.style={to path={|- (\tikztotarget)}}]
%   \matrix[column sep=4mm] {
%     % First row:
%     & & & & & & &  & & & & \node (plus) [terminal] {+};\\
%     % Second row:
%     \node (p1) [point]  {}; &    \node (ui1)   [nonterminal] {unsigned integer}; &
%     \node (p2) [point]  {}; &    \node (dot)   [terminal]    {.};                &
%     \node (p3) [point]  {}; &    \node (digit) [terminal]    {digit};            &
%     \node (p4) [point]  {}; &    \node (p5)    [point]  {};                      &
%     \node (p6) [point]  {}; &    \node (e)     [terminal]    {E};                &
%     \node (p7) [point]  {}; &                                                    &
%     \node (p8) [point]  {}; &    \node (ui2)   [nonterminal] {unsigned integer}; &
%     \node (p9) [point]  {}; &    \node (p10)   [point]       {};\\
%     % Third row:
%     & & & & & & &  & & & & \node (minus)[terminal] {-};\\
%   };
%
%   { [start chain]
%     \chainin (p1);
%     \chainin (ui1)   [join=by tip];
%     \chainin (p2)    [join];
%     \chainin (dot)   [join=by tip];
%     \chainin (p3)    [join];
%     \chainin (digit) [join=by tip];
%     \chainin (p4)    [join];
%     { [start branch=digit loop]
%       \chainin (p3) [join=by {skip loop=-6mm,tip}];
%     }
%     \chainin (p5)    [join,join=with p2 by {skip loop=6mm,tip}];
%     \chainin (p6)    [join];
%     \chainin (e)     [join=by tip];
%     \chainin (p7)    [join];
%     { [start branch=plus]
%       \chainin (plus)  [join=by {vh path,tip}];
%       \chainin (p8)    [join=by {hv path,tip}];
%     }
%     { [start branch=minus]
%       \chainin (minus) [join=by {vh path,tip}];
%       \chainin (p8)    [join=by {hv path,tip}];
%     }
%     \chainin (p8)    [join];
%     \chainin (ui2)   [join=by tip];
%     \chainin (p9)    [join,join=with p6 by {skip loop=-11mm,tip}];
%     \chainin (p10)   [join=by tip];
%   }
% \end{tikzpicture}
% \end{codeexample}
